\documentclass[12pt]{article}
\usepackage[margin=1in]{geometry}
\usepackage{booktabs}
\usepackage{longtable}
\usepackage{float}
\usepackage{setspace}
\setstretch{1.1}
\title{Legends of Noblesse\\AI Tier Matchup Report}
\author{ECE348 Project Simulation Study}
\date{March 31, 2026}
\begin{document}
\maketitle
\section{Introduction}
This report evaluates game balance in \textit{Legends of Noblesse} by comparing three AI skill tiers (bad, mediocre, good) across six requested head-to-head matchups. The central question is whether match outcomes and overall point rates indicate a reasonably fair system rather than a single dominant strategy.
The analysis used \textbf{120} automated games (20 per matchup), under legal engine rules only, with no overtime force logic or deterministic tiebreak. A game ends in a decisive result only when a barracks is successfully breached; otherwise it is recorded as a draw at the configured turn cap.
\section{Rules (Implementation Summary)}
\begin{enumerate}
\item Each player starts with a legal randomized loadout (deck, class, barracks, and 3 battlefields).
\item Every game follows engine phase order: Replenish, Draw, Preparations, Siege, and Field Cleanup.
\item During Draw and Preparations, AIs take only legal actions available in the current phase.
\item During Siege, battalions can only be assigned to legal targets reported by the engine.
\item A decisive winner is declared only by true in-engine barracks capture conditions.
\item No overtime force behavior was used to force attacks or force winners.
\item If no winner is declared by 2000 turns, the game is recorded as a draw for analysis.
\end{enumerate}
\section{Results}
\subsection{Overall Outcomes}
The full 120-game run produced the following aggregate balance metrics:
\begin{itemize}
\item Simulated games: \textbf{120}
\item Non-draw games: \textbf{90}
\item Draws: \textbf{30} (25.00\%)
\item Player 1 decisive win rate (excluding draws): \textbf{26.67\%}
\item Player 2 decisive win rate (excluding draws): \textbf{73.33\%}
\item Player 1 point rate (draw = 0.5): \textbf{32.50\%}
\item Player 2 point rate (draw = 0.5): \textbf{67.50\%}
\item Average game length: \textbf{640.42 turns}
\end{itemize}
\subsection{Matchup Matrix}
\begin{table}[H]
\centering
\caption{20-Game Results per Requested Matchup}
\begin{tabular}{lrrrrrr}
\toprule
Matchup & Games & P1 Wins & P2 Wins & Draws & P1 Win\% & P2 Win\% \\
\midrule
bad vs bad & 20 & 0 & 0 & 20 & 0.0 & 0.0 \\
bad vs mediocre & 20 & 0 & 20 & 0 & 0.0 & 100.0 \\
bad vs good & 20 & 0 & 20 & 0 & 0.0 & 100.0 \\
mediocre vs good & 20 & 6 & 14 & 0 & 30.0 & 70.0 \\
mediocre vs mediocre & 20 & 13 & 7 & 0 & 65.0 & 35.0 \\
good vs good & 20 & 5 & 5 & 10 & 50.0 & 50.0 \\
\bottomrule
\end{tabular}
\end{table}
\subsection{Overall Tier Strength}
\begin{table}[H]
\centering
\caption{Aggregate Tier Performance Across All Appearances}
\begin{tabular}{lrrrrrr}
\toprule
AI Tier & Appearances & Wins & Losses & Draws & Win\% & Point\% \\
\midrule
good & 80 & 44 & 16 & 20 & 73.33 & 67.5 \\
mediocre & 80 & 46 & 34 & 0 & 57.5 & 57.5 \\
bad & 80 & 0 & 40 & 40 & 0.0 & 25.0 \\
\bottomrule
\end{tabular}
\end{table}
\subsection{Overtime Order-Bias Audit (Legacy Forced-Overtime Run)}
To address the question about overtime bias, we also audited the earlier forced-overtime dataset (stored in reports/ai\_tier\_matchups). In that build, overtime began after turn 50, and Player 2 barracks attack resolution occurred before Player 1.
\begin{itemize}
\item Winners in audited run: \textbf{1500}
\item Winners occurring after overtime started: \textbf{832} (55.47\%)
\item Of overtime winners, Player 2 won \textbf{761} vs Player 1 \textbf{71} (91.47\% Player 2 share)
\item At turn 51 (the first overtime turn), \textbf{821} wins were recorded, with Player 2 taking \textbf{750} (91.35\%)
\end{itemize}
These counts quantify how many outcomes were produced in the overtime window where the Player 2 ordering advantage could influence decisive results.
\section{Conclusion}
The strongest tier in this run was good (point rate 67.5\%), while the weakest was bad (point rate 25.0\%). The results show clear skill-tier separation (good $>$ mediocre $>$ bad in aggregate), while preserving legal winner conditions. The analysis therefore answers the assignment prompts with a full rule summary, simulation evidence, and data-backed balance interpretation.
\end{document}
